\section{Exact cofactor occupancy and terminal activation}
\label{sec:tpc59-occupancy}

The critical cutoff makes a cofactor window possible at the desired
exponent.  We now determine exactly what is needed for the literal carrier
to place energy in that window.  This is an orthogonal projection problem
before terminal grouping and therefore admits a sharp rowwise solution.

\subsection{Literal pre-terminal energy}

On one declared source system, let \(\cI_X\) be the fixed retained set of
\((m,r)\)-pairs, \(\cJ_{X,a}\) the fixed orbit class, and \(\cO_X\) the
complete opposite-row cell.  The TPC--53/57 factorization can be written
\begin{equation}
 [G_X^L]_{m,r,j,\gamma}
 =c_{m,j,\gamma}\zeta_m,
 \qquad
 (m,r)\in\cI_X,\quad j\in\cJ_{X,a},\quad\gamma\in\cO_X,
 \label{eq:tpc59-literal-preterminal-factorization}
\end{equation}
where \(c_{m,j,\gamma}\) contains the source shape with its M\"obius and
prime-logarithm factors, literal mask, opposite-row factor, orbit phase, and
four-factor profile, but not the bounded row coefficient \(\zeta_m\).  The
essential fact is that \(c_{m,j,\gamma}\) has no \(r\)-dependence
\citep{WangTPC53,WangTPC57}.

Put
\begin{equation}
 B_m:=\sum_{j\in\cJ_{X,a}}\sum_{\gamma\in\cO_X}
       |c_{m,j,\gamma}|^2,
 \qquad
 N(m):=\#\{r:(m,r)\in\cI_X\}.
 \label{eq:tpc59-Bm-Nm}
\end{equation}
Then the raw parent energy is exactly
\begin{equation}
 \boxed{
 \cE_{\rm par}(\zeta)
 =\sum_mB_m|\zeta_m|^2N(m).}
 \label{eq:tpc59-parent-energy}
\end{equation}
Repeated cofactors are distinct orthogonal atoms in this identity.

For any fixed cofactor set \(\cR\subset(0,\infty)\), define
\begin{equation}
 N_{\cR}(m)
 :=\#\{r:(m,r)\in\cI_X,\ r\in\cR\}
 \label{eq:tpc59-compatible-row-multiplicity}
\end{equation}
and the projected energy
\begin{equation}
 \boxed{
 \cE_{\cR}(\zeta)
 =\sum_mB_m|\zeta_m|^2N_{\cR}(m).}
 \label{eq:tpc59-compatible-energy}
\end{equation}
The set \(\cR\) may be the safe critical window, a union of dyadic blocks,
or one distinguished block.  It is fixed independently of \(\zeta\).

\begin{theorem}[Sharp rowwise occupancy minimax]
\label{thm:tpc59-rowwise-occupancy-minimax}
Assume \(\cE_{\rm par}\) is not identically zero on the coefficient cube
\(|\zeta_m|\le1\).  The optimal coefficient-uniform projection constant is
\begin{equation}
 \boxed{
 \begin{aligned}
 \beta_{\cR}
 &:={}
 \inf_{\substack{|\zeta_m|\le1\\\cE_{\rm par}(\zeta)>0}}
 \frac{\cE_{\cR}(\zeta)}{\cE_{\rm par}(\zeta)}\\
 &=\min_{m:B_mN(m)>0}\frac{N_{\cR}(m)}{N(m)}.
 \end{aligned}}
 \label{eq:tpc59-beta-R-minimax}
\end{equation}
In particular,
\begin{equation}
 \beta_{\cR}>0
 \quad\Longleftrightarrow\quad
 N_{\cR}(m)>0
 \quad\text{for every row with }B_mN(m)>0.
 \label{eq:tpc59-beta-positive-test}
\end{equation}
The formula is sharp at every finite \(X\).
\end{theorem}

\begin{proof}
For \(B_mN(m)>0\), write
\[
 \rho_m:=\frac{N_{\cR}(m)}{N(m)},
 \qquad
 w_m:=B_mN(m)|\zeta_m|^2.
\]
Whenever the denominator is positive,
\begin{equation}
 \frac{\cE_{\cR}(\zeta)}{\cE_{\rm par}(\zeta)}
 =\frac{\sum_mw_m\rho_m}{\sum_mw_m},
 \label{eq:tpc59-occupancy-convex-combination}
\end{equation}
which is a convex combination of the row ratios.  It is at least their
minimum.  Conversely, choose a minimizing row \(m_*\), set
\(\zeta_{m_*}=1\), and set every other row coefficient to zero.  Then the
ratio equals \(\rho_{m_*}\).  The positivity test follows immediately.
\end{proof}

For one distinguished coefficient vector, define the exact weaker ratio
\begin{equation}
 \rho_{\cR}(\zeta)
 :=\frac{\sum_mB_m|\zeta_m|^2N_{\cR}(m)}
         {\sum_mB_m|\zeta_m|^2N(m)},
 \label{eq:tpc59-distinguished-occupancy}
\end{equation}
whenever the denominator is positive.  Then
\begin{equation}
 \beta_{\cR}\le\rho_{\cR}(\zeta)\le1.
 \label{eq:tpc59-uniform-distinguished-order}
\end{equation}
A favorable distinguished ratio does not prove a uniform carrier theorem.

\begin{corollary}[Sharp zero-occupancy obstruction]
\label{cor:tpc59-zero-occupancy-obstruction}
If one active row has retained cofactors but none in \(\cR\), then
\(\beta_{\cR}=0\).  Hence geometric compatibility of \(\cR\), positive
cardinality of \(\cR\) in other rows, and boundedness of \(\zeta\) cannot
imply a positive uniform parent-to-window transfer.
\end{corollary}

This is why the word ``relocation'' must be handled carefully.  The map is
the coordinate projection onto already existing \(r\)-labels; it does not
send an atom at one cofactor to another cofactor.  A fixed-factor cutoff
retuning can change which existing labels are compatible, but it cannot
create their energy.

\subsection{Terminal activation as a second rowwise minimax}

For a pre-terminal coordinate \(u=(m,r,j,\gamma)\), put
\begin{equation}
 n_u:=mj+h_0,
 \qquad
 p_u:=\frac{n_u}{r}
 \label{eq:tpc59-terminal-candidate}
\end{equation}
when the quotient is integral.  Let \(\one_{\rm term}^{>y}(u)\) be the
indicator that \(u\) satisfies the complete inherited terminal support:
\(p_u\) is prime, \(p_u>y\), \(\Lambda_X(n_u)\ne0\), and all required
squarefree, coprimality, and support conditions hold.  Define the rowwise
terminal weight inside \(\cR\) by
\begin{equation}
 T_{\cR}(m)
 :=\sum_{\substack{r:(m,r)\in\cI_X\\r\in\cR}}
   \sum_{j,\gamma}
   \one_{\rm term}^{>y}(m,r,j,\gamma)
   |\Lambda_X(n_u)|^2|c_{m,j,\gamma}|^2.
 \label{eq:tpc59-row-terminal-weight}
\end{equation}
The literal terminal amplitude is
\begin{equation}
 x_u=\mu(r)\Lambda_X(n_u)c_{m,j,\gamma}\zeta_m,
 \label{eq:tpc59-terminal-amplitude}
\end{equation}
so the raw terminal diagonal in \(\cR\) is
\begin{equation}
 \boxed{
 \cD_{\cR}^{\rm term}(\zeta)
 =\sum_mT_{\cR}(m)|\zeta_m|^2.}
 \label{eq:tpc59-terminal-energy-row-sum}
\end{equation}

\begin{proposition}[Exact terminal-activation minimax]
\label{prop:tpc59-terminal-activation-minimax}
Assume \(\cE_{\cR}\) is not identically zero.  The optimal uniform lower
constant from compatible pre-terminal energy to its very-high terminal
diagonal is
\begin{equation}
 \boxed{
 \tau_{\cR}
 :=\inf_{\substack{|\zeta_m|\le1\\\cE_{\cR}(\zeta)>0}}
   \frac{\cD_{\cR}^{\rm term}(\zeta)}{\cE_{\cR}(\zeta)}
 =\min_{m:B_mN_{\cR}(m)>0}
   \frac{T_{\cR}(m)}{B_mN_{\cR}(m)}.}
 \label{eq:tpc59-tau-R-minimax}
\end{equation}
The exact combined parent-to-terminal constant is
\begin{equation}
 \boxed{
 \chi_{\cR}
 :=\inf_{\substack{|\zeta_m|\le1\\\cE_{\rm par}(\zeta)>0}}
   \frac{\cD_{\cR}^{\rm term}(\zeta)}{\cE_{\rm par}(\zeta)}
 =\min_{m:B_mN(m)>0}\frac{T_{\cR}(m)}{B_mN(m)},}
 \label{eq:tpc59-combined-minimax}
\end{equation}
with the convention \(T_{\cR}(m)=0\) if
\(N_{\cR}(m)=0\).  Moreover,
\begin{equation}
 \chi_{\cR}\ge\beta_{\cR}\tau_{\cR}.
 \label{eq:tpc59-combined-product-lower}
\end{equation}
None of these exact formulas asserts that the displayed constant is
positive.
\end{proposition}

\begin{proof}
The two minimax formulas follow from the same weighted-average argument as
\cref{thm:tpc59-rowwise-occupancy-minimax}, concentrating \(\zeta\) on one
row for sharpness.  If an active parent row has \(N_{\cR}(m)=0\), then
\(\beta_{\cR}=\chi_{\cR}=0\), and
\eqref{eq:tpc59-combined-product-lower} is immediate.  Otherwise every
active parent row has \(N_{\cR}(m)>0\), and
\[
 \frac{T_{\cR}(m)}{B_mN(m)}
 =\frac{N_{\cR}(m)}{N(m)}
  \frac{T_{\cR}(m)}{B_mN_{\cR}(m)}.
\]
Taking rowwise minima proves
\eqref{eq:tpc59-combined-product-lower}.
\end{proof}

The inherited hypotheses impose no prime-distribution lower bound on
\(T_{\cR}(m)\).  Thus \(\tau_{\cR}\) and \(\chi_{\cR}\) may vanish even
when \(\beta_{\cR}>0\).  Conversely, a positive terminal face in selected
rows does not repair a row with zero compatible occupancy.  This formally
separates the last two gates in \eqref{eq:tpc59-three-gates}.

\begin{remark}[Activation constants are not probabilities]
The literal multiplier \(|\Lambda_X(n)|^2\) remains inside
\(T_{\cR}(m)\).  Consequently \(\tau_{\cR}\) and \(\chi_{\cR}\) are
energy-transfer constants and can exceed one; only their possible vanishing
and their polynomial lower scale are relevant to the loss ledger.
\end{remark}
